\section{Introduction}

Expectation thresholds measure how cheaply an increasing family can be
covered by elementary containment events. Their fractional relaxation
allows a member to be covered by a total weight of one, distributed
among its subsets. Talagrand conjectured that this relaxation changes
the expectation threshold by at most a universal factor
\cite[Conjecture~6.3]{Talagrand2010}. We prove this conjecture.

Let $X$ be a finite nonempty set and let $\F$ be a nonempty proper
increasing subfamily of $2^X$. For a family $\G\subseteq2^X$ and a
function $g:2^X\to[0,1]$, define
\[
 \langle\G\rangle
 =\{H\subseteq X:\exists S\in\G,\ S\subseteq H\},
 \qquad
 c_p(\G)=\sum_{S\in\G}p^{|S|},
\]
and
\[
 \langle g\rangle
 =\left\{H\subseteq X:\sum_{S\subseteq H}g(S)\ge1\right\},
 \qquad
 w(g,p)=\sum_{S\subseteq X}g(S)p^{|S|}.
\]
Throughout, $p\in[0,1]$ and $p^0=1$, including when $p=0$.
Thus a cover supplies, inside each member of $\F$, at least one
generator $S\in\G$. If $X_p$ includes the elements of $X$ independently
with probability $p$, the cost of that generator is
$\mathbb P(S\subseteq X_p)=p^{|S|}$.
We say that $\F$ is \emph{$p$-small} if there is a family $\G$ with
$\F\subseteq\langle\G\rangle$ and $c_p(\G)\le1/2$.
Its integral and fractional expectation thresholds are
\begin{align*}
 q(\F)
 &=\max\{p\in[0,1]:\F\text{ is }p\text{-small}\},\\
 q_f(\F)
 &=\max\{p\in[0,1]:\exists g:2^X\to[0,1],
       \F\subseteq\langle g\rangle,
       \ w(g,p)\le1/2\}.
\end{align*}
In particular, $q(\F)\le q_f(\F)$, by taking $g$ to be the indicator
of an integral cover.

These maxima are well-defined. Since $\varnothing\notin\F$, both
feasible parameter sets contain zero. Integral feasibility is a finite
union of closed subsets of $[0,1]$. Fractional feasibility is the
projection of the compact set of pairs $(p,g)$ satisfying the displayed
constraints in $[0,1]\times[0,1]^{2^X}$.

\Needspace{9\baselineskip}
\begin{theorem}\label{thm:main}
For every finite nonempty $X$ and every nonempty proper increasing
family $\F\subseteq2^X$,
\[
 q_f(\F)\le25\cdot512^4\,q(\F).
\]
Both thresholds use the budget $1/2$. The constant is independent of
$|X|$, the sizes of the minimal members of $\F$, and the support of
the fractional cover.
\end{theorem}

The conclusion preserves the covering budget: every fractional cover of
cost at most $1/2$ at a parameter $r$ admits an integral cover of cost at
most $1/2$ at $r/(25\cdot512^4)$.

\subsection{Thresholds, selectors, and earlier comparisons}

Expectation thresholds arose in the study of when a random set first
has a prescribed increasing property. Let $p_c(\F)$ denote the unique
parameter for which $\mathbb P(X_p\in\F)=1/2$. A union bound applied to
an integral cover gives $q(\F)\le p_c(\F)$; taking expectations in a
fractional cover gives $q_f(\F)\le p_c(\F)$ as well.
Kahn and Kalai~\cite{KahnKalai2007} conjectured the upper bound
$p_c(\F)=O(q(\F)\log(2|X|))$.
Frankston, Kahn, Narayanan, and Park~\cite{FKNP2021}, building on the
sunflower methods of Alweiss, Lovett, Wu, and Zhang~\cite{ALWZ2021},
proved the fractional bound
\[
 p_c(\F)=O(q_f(\F)\log\ell),
 \qquad
 \ell=\max\{2,\max\{|H|:H\text{ is a minimal member of }\F\}\}.
\]
Park and Pham~\cite{ParkPham2024threshold} proved the corresponding
bound with $q$ in place of $q_f$, resolving the Kahn--Kalai conjecture.
These theorems estimate the random-set threshold. Talagrand's comparison
asks whether fractional certificates themselves admit integral
replacements at a constant loss.

Early results obtained such replacements by exploiting the structure
of the sets with positive fractional weight. Talagrand settled the case
of singleton sets~\cite[Section~6]{Talagrand2010}; DeMarco and
Kahn~\cite{DeMarcoKahn2015} treated clique-counting certificates, and
Frankston, Kahn, and Park~\cite{FrankstonKahnPark2022} treated arbitrary
weights on pairs. Fischer and Person extended these methods to nearly
linear uniform hypergraphs~\cite{FischerPerson2025} and further clique
cases~\cite{FischerPerson2025further}. They also proved the comparison
with high probability for constant weights on random uniform supports
\cite{FischerPerson2026random}.

Selector processes supply a different route: they seek a member whose
assigned vertex mass is substantially captured by a random subset.
Park and Pham proved Talagrand's selector-process conjecture
\cite{ParkPham2024selector}, and Bednorz, Martynek, and
Meller~\cite{BMM2022} gave a proof using truncated weights.
Dubroff, Kahn, and Park~\cite[Theorem~1.2]{DKP2024} used selector
estimates to round any fractional cover whose positive coefficients
occur only on sets of cardinality at most $t$, with a loss depending on
$t$. Pham sharpened this loss to $O(\log(2t))$
\cite[arXiv version~1, Theorem~1.2]{Pham2024}.
Here $t$ bounds the cardinalities of the weighted sets, not their number
or the sizes of the minimal members of $\F$. For example, for
$1\le k\le|X|$, the family $\{H\subseteq X:|H|\ge k\}$ has the
fractional cover $g(\{x\})=1/k$, zero elsewhere: its weighted sets are
singletons, while its minimal members have cardinality $k$.

Fischer and Person's sampling argument for sufficiently large sets
\cite[Theorem~10]{FischerPerson2025}, combined with Pham's logarithmic
rounding bound, gives a loss of order $\log\log(4|X|)$. This consequence
is recorded in the expanded version of Pham's paper
\cite[Corollary~1.8]{Pham2026}. More recently, Park
\cite[Theorem~1.2]{Park2026} obtained the dimension-independent bound
\[
 q_f(\F)\le Kq(\F)\max\{1,\log\log(1/q(\F))\}
\]
for a universal $K$. The factor here still grows as $q(\F)$ tends to
zero. Theorem~\ref{thm:main} removes both kinds of dependence.
Section~\ref{sec:consequence} applies it to Li's fractional
discrete-convexity theorem~\cite{Li2026}: for any class
$\mathcal A\subseteq2^X$ with $\P(X_p\in\mathcal A)>1/2$, the sets
not contained in a union of two members of $\mathcal A$ have a small
integral cover at a universal constant multiple of $p$.

\begin{remark}[Talagrand's original formulation]
For $p>0$, the weak-smallness probability weights $\beta_S$ in
\cite[Definition~6.1]{Talagrand2010} give fractional coefficients
$\beta_S/(2p^{|S|})$; capping these at one preserves coverage and does
not increase cost. Conversely, the masses $2g(S)p^{|S|}$ can be
completed to a probability distribution by adding the deficit at the
empty set. Passing from an arbitrary class to its upward closure
preserves both cover constraints. Thus the increasing-family formulation
of Theorem~\ref{thm:main} also resolves Conjecture~6.3 in its original
form.
\end{remark}

\subsection{The selector estimate and the rounding mechanism}

Our selector estimate controls one member simultaneously at several
scales. Assign to each $H\in\F$ a probability mass $\lambda_H$ supported
on $H$. Independently color the vertices by integers
$a(x)\in\{1,\ldots,s+1\}$, and let
$A_i=\{x:a(x)\le i\}$. Theorem~\ref{thm:multiscale-selector} specifies
geometrically increasing color probabilities for which failure of
$p$-smallness forces, with probability at least $9/10$, a single $H$
satisfying
\[
 \lambda_H(A_i)\ge1-2^{-i}\quad(1\le i\le s).
\]
The uncaptured masses are summable, so this member has bounded mean
color:
\[
 \sum_x a(x)\lambda_H(x)
 =1+\sum_{i=1}^s\bigl(1-\lambda_H(A_i)\bigr)\le2.
\]
Separate estimates at individual scales could choose different members
and would not give this bound for any one member.

To apply the estimate to a fractional cover of cost at most $1/2$ at
parameter $r$, distribute each nonempty
weight $g(S)$ equally among the vertices of $S$ and normalize on each
member, following Pham~\cite[arXiv version~1, Section~3]{Pham2024}.
The mean vertex color then equals the weighted mean of the average
colors of those sets. Thus a fixed positive amount of fractional weight
lies on sets with average color at most four. For a suitable fixed base $B>1$,
the independent variables $Y_x=B^{4-a(x)}$ satisfy
$\prod_{x\in S}Y_x\ge1$ on all these sets. Their common mean can at the
same time be made a constant multiple of $p$ by a suitable choice of
the number of colors. Taking $p$ to be a sufficiently small constant
multiple of $r$, independence bounds the expectation of
\[
 Z=\sum_{\varnothing\ne S\subseteq X}g(S)\prod_{x\in S}Y_x
\]
from above using the fractional covering cost, while the selector event
bounds it from below. Section~\ref{sec:rounding} makes these bounds
incompatible. Only nonemptiness of the weighted sets enters this final
comparison, so their cardinalities need no upper bound.

The proof of the selector estimate in Section~\ref{sec:selector}
develops the minimum-fragment method of Park and Pham
\cite{ParkPham2024selector} and Pham's towers of fragments
\cite[arXiv version~1, Section~2]{Pham2024}, together with the maximal
truncations of Bednorz, Martynek, and Meller~\cite[Section~3.1]{BMM2022}.
For each coloring and each member, we move vertices to earlier colors
until a set of truncated inequalities holds, minimizing the total
number of color levels moved. The resulting changed sets form a cover.
To estimate its cost, we record the final coloring and the numbers of
vertices making each possible color change. These data determine one
common member satisfying the truncated inequalities for the entire
group of original colorings.
Maximal truncation bounds the number of candidate vertices at each color,
and a weighted
AM--GM estimate pays for their allocation among the scales. Keeping all
color changes in the same record is what permits the simultaneous
estimate.

The random vertex colors used below are independent.
We write $\lambda_H(A)=\sum_{x\in A}\lambda_H(x)$.
Empty sums are zero and empty products are one.
